
\subsubsection{6.1 Interpretation of the PARTIAL Gate Result}

The h-e1 gate result is PARTIAL: infrastructure validated, behavioral gate not met. The root cause is a single missing prerequisite: the SFT checkpoint specified as a controlled variable in the hypothesis design was not available at experiment time. The mechanism is not falsified. The prescreening protocol produces the correct output (S\_term = 0 for all problems with a model that passes zero tests), and the gate failure correctly reflects the absence of qualifying problems.

The PARTIAL designation differs from a FAIL in the following respect: FAIL would mean the mechanism does not work as expected (e.g., S\_term is computed incorrectly, or the variance ratio is not elevated even for problems in the target window). PARTIAL means the infrastructure and mechanism are correct, but a prerequisite condition for testing the central claim is not satisfied. The distinction determines the recommended next action: FAIL routes to fundamental hypothesis revision, while PARTIAL routes to prerequisite resolution.

\subsubsection{6.2 The SFT Prerequisite}

The 0\% pass rate finding has two probable contributing factors. First, format incompatibility: the instruction-tuned model generates responses in a chat format that the execution sandbox does not parse. Second, capability gap: even with correct formatting, the model may not be able to solve APPS introductory problems from the instruction-tuned base. SFT on APPS-format solutions addresses both factors simultaneously by training the model to produce raw executable Python in the APPS output format.

The finding is directionally consistent with Du et al. [2025], who report that SFT initialization is required before GRPO on competitive programming problems produces meaningful signal. The present study cannot separately attribute the 0\% result to format incompatibility versus capability, as this would require a format diagnostic experiment (e.g., extracting code blocks from model outputs before execution). This diagnostic is identified as a recommended next step.

The practical implication is that researchers applying GRPO to APPS or similar competitive programming benchmarks with instruction-tuned base models should not assume the base model is ready for execution-based reward evaluation without verifying format compatibility.

\subsubsection{6.3 Scope of Completed Contributions}

Two of the four stated contributions are fully substantiated by the evidence in this study:

\textbf{Contribution 2 (Binomial variance proof):} The analytical derivation in Section 3.1 is mathematically complete. The variance ratio formula $\rho$(q, T) = q($1-q)$ / [T $\cdot q^T$($1-q^{T})]$ follows from the definitions of R\_ratio and R\_binary and the properties of the Binomial distribution. This derivation does not require empirical confirmation.

\textbf{Contribution 3 (Production-ready prescreening infrastructure):} The pipeline validation is confirmed by 15/15 task completion and 67/67 test passes.

Two contributions are substantiated in part:

\textbf{Contribution 1 (S\_term formalization):} The formal definition and the target window are analytically motivated. Whether S\_term $\in$ [0.3, 0.55] is achievable for a sufficient number of APPS introductory problems with an SFT-initialized model is empirically unverified.

\textbf{Contribution 4 (SFT prerequisite discovery):} The 0\% pass rate is confirmed empirically. The attribution of this result to format incompatibility versus capability is not confirmed and remains a probable interpretation rather than an established fact.

\subsubsection{6.4 Limitations}

\textbf{L1 --- SFT checkpoint gap.} All five sub-hypotheses are blocked until an SFT checkpoint on APPS introductory solutions is available. The checkpoint was a specified controlled variable whose generation was not included as a pipeline task.

\textbf{L2 --- Format incompatibility unconfirmed.} The 0\% pass rate probably reflects output format mismatch, but this has not been verified through a targeted diagnostic. A format diagnostic (code block extraction prior to execution) is needed to distinguish format failure from capability failure.

\textbf{L3 --- Tractability window assumption unverified.} The window S\_term $\in$ [0.3, 0.55] is analytically derived but not empirically confirmed to be achievable for a meaningful fraction of APPS introductory problems with an SFT-initialized 7B model.

\textbf{L4 --- Limited problem sample.} 300 of 1,923 available problems were processed (15.6\%). After SFT enables non-zero pass rates, gate metric estimates from 300 problems may have high standard error if qualifying problems are sparse.

\textbf{L5 --- Single seed for prescreening.} The prescreening experiment used seed=42 only. Sampling variance in the prescreening result is unquantified.

\textbf{L6 --- Execution timeout.} The 5-second per-test-case timeout may cause correct but computationally intensive solutions to be counted as failures, potentially underestimating S\_term for problems with slow solutions.

\textbf{L7 --- Single model family.} All experiments use Qwen2.5-Coder-7B-Instruct. The Binomial variance mechanism is model-agnostic, but the tractability distribution and optimal prescreening window will vary across models and sizes.

